
This work investigated whether signals related to human agency preservation can be extracted from existing AI preference data. The Bidirectional Alignment Index achieves high extraction reliability (mean AUROC 0.98) and survives adversarial probing designed to remove reward-predictive variance (AUROC 0.99 post-gradient-reversal). However, semantic analysis reveals that high-BAI responses cluster by generic conversational patterns rather than interpretable agency vocabulary (0\% agency pattern rate).

This combination of findings refines understanding of what surface-level proxies can and cannot accomplish: variance orthogonal to reward can be reliably extracted, but this variance may not capture the intended construct. The gap between extractability and interpretability is the core lesson.

The methodological contributions---adversarial probing for independence verification and semantic clustering for content validation---remain applicable to future bidirectional alignment research. The negative result on semantic coherence directs future work toward richer operationalizations: embedding-based classifiers, LLM-as-judge evaluations, or human annotation of functional agency.
